\textbf{Phần 1. Trắc nghiệm lựa chọn}
\setcounterpageref{ex}{0}

\begin{ex} %[5D10N58]
Kết quả của phép tính $1$ giờ $20$ phút $\times 3$ là
\choice
{$3$ giờ $20$ phút}
{$3$ giờ $60$ phút}
{$4$ giờ $20$ phút}
{\True $4$ giờ}
\loigiai{
$1$ giờ $20$ phút $\times 3 = 3$ giờ $60$ phút $= 4$ giờ.\\
\textbf{Đáp án của câu hỏi là D}
}
\end{ex}

\begin{ex} %[5D10N58]
Kết quả của phép tính $8$ phút $40$ giây $: 4$ là
\choice
{\True $2$ phút $10$ giây}
{$2$ phút}
{$2$ phút $20$ giây}
{$1$ phút $10$ giây}
\loigiai{
$8$ phút $40$ giây $: 4 = 2$ phút $10$ giây.\\
\textbf{Đáp án của câu hỏi là A}
}
\end{ex}

\begin{ex} %[5D10H58]
Kết quả của phép tính $2$ giờ $45$ phút $\times 4$ là
\choice
{$8$ giờ $180$ phút}
{$10$ giờ $45$ phút}
{\True $11$ giờ}
{$12$ giờ}
\loigiai{
$2$ giờ $45$ phút $\times 4 = 8$ giờ $180$ phút $= 11$ giờ.\\
\textbf{Đáp án của câu hỏi là C}
}
\end{ex}

\begin{ex} %[5D10H58]
Kết quả của phép tính $7$ giờ $30$ phút $: 5$ là
\choice
{$1$ giờ $6$ phút}
{$1$ giờ $15$ phút}
{$1$ giờ $45$ phút}
{\True $1$ giờ $30$ phút}
\loigiai{
$7$ giờ $: 5 = 1$ giờ (dư 2 giờ $= 120$ phút).\\
$(120 + 30) : 5 = 150 : 5 = 30$ phút.\\
Kết quả: $1$ giờ $30$ phút.\\
\textbf{Đáp án của câu hỏi là D}
}
\end{ex}

\begin{ex} %[5D10H58]
Một người thợ làm một sản phẩm hết $45$ phút. Hỏi người đó làm $6$ sản phẩm như thế hết bao nhiêu thời gian?
\choice
{$4$ giờ}
{\True $4$ giờ $30$ phút}
{$4$ giờ $15$ phút}
{$5$ giờ}
\loigiai{
$45 \times 6 = 270$ phút $= 4$ giờ $30$ phút.\\
\textbf{Đáp án của câu hỏi là B}
}
\end{ex}

\begin{ex} %[5D10H58]
Một vận động viên hoàn thành $4$ vòng chạy quanh sân vận động hết $9$ phút $20$ giây. Trung bình mỗi vòng vận động viên đó chạy hết bao lâu?
\choice
{$2$ phút $15$ giây}
{$2$ phút $30$ giây}
{\True $2$ phút $20$ giây}
{$2$ phút $10$ giây}
\loigiai{
$9$ phút $20$ giây $: 4 = 560$ giây $: 4 = 140$ giây $= 2$ phút $20$ giây.\\
\textbf{Đáp án của câu hỏi là C}
}
\end{ex}

\begin{ex} %[5D10V58]
Một vòi nước chảy vào bể cứ $15$ phút thì được $120\text{ l}$ nước. Hỏi vòi đó chảy trong $1$ giờ $15$ phút thì được bao nhiêu lít nước?
\choice
{$480\text{ l}$}
{\True $600\text{ l}$}
{$720\text{ l}$}
{$500\text{ l}$}
\loigiai{
$1$ giờ $15$ phút $= 75$ phút.\\
$75$ phút gấp $15$ phút số lần là: $75 : 15 = 5$ (lần).\\
Số lít nước chảy được: $120 \times 5 = 600\text{ (l)}$.\\
\textbf{Đáp án của câu hỏi là B}
}
\end{ex}

\begin{ex} %[5D10V58]
Bác Ba làm $5$ chiếc ghế hết tất cả $8$ giờ $20$ phút. Bác Tư làm $6$ chiếc ghế như thế hết $9$ giờ $30$ phút. So sánh thời gian làm $1$ chiếc ghế của hai bác:
\choice
{Bác Ba làm nhanh hơn bác Tư}
{Hai bác làm nhanh như nhau}
{\True Bác Tư làm nhanh hơn bác Ba}
{Không so sánh được}
\loigiai{
Bác Ba làm 1 ghế: $8$ giờ $20$ phút $: 5 = 500$ phút $: 5 = 100$ phút $= 1$ giờ $40$ phút.\\
Bác Tư làm 1 ghế: $9$ giờ $30$ phút $: 6 = 570$ phút $: 6 = 95$ phút $= 1$ giờ $35$ phút.\\
Vì $95$ phút $< 100$ phút nên bác Tư làm nhanh hơn bác Ba.\\
\textbf{Đáp án của câu hỏi là C}
}
\end{ex}

\begin{ex} %[5D10V58]
Một đội công nhân lắp ráp $8$ bộ máy hết $14$ giờ. Hỏi đội đó lắp ráp $12$ bộ máy như thế hết bao nhiêu thời gian?
\choice
{\True $21$ giờ}
{$20$ giờ}
{$22$ giờ}
{$24$ giờ}
\loigiai{
Lắp 1 bộ máy hết: $14 : 8 = 1{,}75$ giờ.\\
Lắp 12 bộ máy hết: $1{,}75 \times 12 = 21$ giờ.\\
\textbf{Đáp án của câu hỏi là A}
}
\end{ex}

\begin{ex} %[5D10V58]
Tính giá trị của biểu thức: $(2$ giờ $15$ phút $+ 1$ giờ $45$ phút$) \times 3$.
\choice
{$11$ giờ}
{$10$ giờ $30$ phút}
{$12$ giờ $30$ phút}
{\True $12$ giờ}
\loigiai{
$2$ giờ $15$ phút $+ 1$ giờ $45$ phút $= 4$ giờ.\\
$4$ giờ $\times 3 = 12$ giờ.\\
\textbf{Đáp án của câu hỏi là D}
}
\end{ex}

\begin{ex} %[5D10V58]
Tính giá trị của biểu thức: $(15$ giờ $20$ phút $- 3$ giờ $40$ phút$) : 5$.
\choice
{$2$ giờ $15$ phút}
{\True $2$ giờ $20$ phút}
{$2$ giờ $25$ phút}
{$2$ giờ $30$ phút}
\loigiai{
$15$ giờ $20$ phút $- 3$ giờ $40$ phút $= 11$ giờ $40$ phút $= 700$ phút.\\
$700$ phút $: 5 = 140$ phút $= 2$ giờ $20$ phút.\\
\textbf{Đáp án của câu hỏi là B}
}
\end{ex}

\begin{ex} %[5D10C58]
Một người thợ đóng $3$ cái bàn và $5$ cái ghế hết tất cả $17$ giờ $30$ phút. Biết thời gian đóng $1$ cái bàn gấp đôi thời gian đóng $1$ cái ghế. Thời gian người thợ đó đóng $1$ cái bàn là
\choice
{\True $3$ giờ $10$ phút}
{$1$ giờ $35$ phút}
{$2$ giờ $45$ phút}
{$3$ giờ}
\loigiai{
Thời gian đóng $3$ cái bàn bằng thời gian đóng $3 \times 2 = 6$ cái ghế.\\
Tổng thời gian tương đương đóng $6 + 5 = 11$ cái ghế.\\
Thời gian đóng 1 cái ghế: $17$ giờ $30$ phút (nếu đề là 17 giờ 25 phút $= 1045$ phút $: 11 = 95$ phút $\Rightarrow$ bàn $= 190$ phút $= 3$ giờ $10$ phút).\\
\textbf{Đáp án của câu hỏi là A}
}
\end{ex}

\textbf{Phần 2. Trắc nghiệm đúng sai}
\setcounterpageref{ex}{0}

\begin{ex} %[5D10H58]
Xét tính đúng sai của các phép nhân, chia số đo thời gian:
\choiceTF
{\True $1$ giờ $25$ phút $\times 4 = 5$ giờ $40$ phút.}
{$14$ phút $30$ giây $: 3 = 4$ phút $10$ giây.}
{\True $14$ phút $30$ giây $: 3 = 4$ phút $50$ giây.}
{\True $3$ ngày $12$ giờ $\times 2 = 7$ ngày.}
\loigiai{
\begin{itemize}
    \item $1$h $25$p $\times 4 = 4$h $100$p $= 5$h $40$p. Đúng.
    \item $14$p $30$s $= 870$s $: 3 = 290$s $= 4$p $50$s $\neq 4$p $10$s. Ý 2 Sai, Ý 3 Đúng.
    \item $3$ ngày $12$h $\times 2 = 6$ ngày $24$h $= 7$ ngày. Đúng.
\end{itemize}
\begin{itemize}
    \item a) \textbf{Đúng}.
    \item b) \textbf{Sai}.
    \item c) \textbf{Đúng}.
    \item d) \textbf{Đúng}.
\end{itemize}
}
\end{ex}

\begin{ex} %[5D10V58]
Một người may $4$ chiếc áo cùng loại hết $6$ giờ. Xét tính đúng sai:
\choiceTF
{\True Trung bình may $1$ chiếc áo hết $1{,}5$ giờ $= 1$ giờ $30$ phút.}
{Nếu may $6$ chiếc áo như thế thì hết $8$ giờ.}
{\True Nếu may $6$ chiếc áo như thế thì hết $9$ giờ.}
{\True Nếu may trong $12$ giờ thì được $8$ chiếc áo.}
\loigiai{
\begin{itemize}
    \item 1 áo: $6 : 4 = 1{,}5$ giờ $= 1$ giờ $30$ phút. Đúng.
    \item 6 áo: $1{,}5 \times 6 = 9$ giờ $\neq 8$ giờ. Ý 2 Sai, Ý 3 Đúng.
    \item Trong 12 giờ: $12 : 1{,}5 = 8$ chiếc áo. Đúng.
\end{itemize}
\begin{itemize}
    \item a) \textbf{Đúng}.
    \item b) \textbf{Sai}.
    \item c) \textbf{Đúng}.
    \item d) \textbf{Đúng}.
\end{itemize}
}
\end{ex}

\textbf{Phần 3. Trả lời ngắn}
\setcounterpageref{ex}{0}

\begin{ex} %[5D10H58]
Tính kết quả phép tính: $2$ giờ $36$ phút $\times 5$ theo đơn vị giờ.
\shortans{13}
\loigiai{
\textbf{Đáp số:} $13$\\
$2{,}6 \times 5 = 13$ giờ.
}
\end{ex}

\begin{ex} %[5D10H58]
Tính kết quả phép tính: $16$ giờ $15$ phút $: 5$ theo phút.
\shortans{195}
\loigiai{
\textbf{Đáp số:} $195$\\
$16 \times 60 + 15 = 975$ phút $\Rightarrow 975 : 5 = 195$ phút.
}
\end{ex}

\begin{ex} %[5D10V58]
Một máy đóng gói tự động đóng được $1$ gói bánh mất $4$ giây. Hỏi trong $15$ phút máy đó đóng được bao nhiêu gói bánh?
\shortans{225}
\loigiai{
\textbf{Đáp số:} $225$\\
$15$ phút $= 900$ giây $\Rightarrow 900 : 4 = 225$ gói bánh.
}
\end{ex}

\begin{ex} %[5D10V58]
Một người làm $3$ dụng cụ hết $5$ giờ $15$ phút. Hỏi làm $8$ dụng cụ như thế hết bao nhiêu giờ?
\shortans{14}
\loigiai{
\textbf{Đáp số:} $14$\\
1 dụng cụ: $5{,}25 : 3 = 1{,}75$ giờ $\Rightarrow 8$ dụng cụ: $1{,}75 \times 8 = 14$ giờ.
}
\end{ex}

\begin{ex} %[5D10V58]
Một chiếc xe máy chạy $3$ vòng trường đua hết $7$ phút $30$ giây. Hỏi xe máy đó chạy $10$ vòng như thế hết bao nhiêu phút?
\shortans{25}
\loigiai{
\textbf{Đáp số:} $25$\\
1 vòng: $7{,}5 : 3 = 2{,}5$ phút $\Rightarrow 10$ vòng: $2{,}5 \times 10 = 25$ phút.
}
\end{ex}

\begin{ex} %[5D10C58]
Một nhóm thợ gồm $4$ người cùng làm xong một công việc hết $6$ giờ $45$ phút. Nếu nhóm có $6$ người làm với cùng năng suất như nhau thì hoàn thành công việc đó hết bao nhiêu phút?
\shortans{270}
\loigiai{
\textbf{Đáp số:} $270$\\
Tổng thời gian 1 người cần: $405 \times 4 = 1\,620$ phút.\\
Thời gian 6 người làm: $1\,620 : 6 = 270$ phút.
}
\end{ex}
